\label{sec:introduction}

Congressional apportionment is a constitutional requirement that runs on a decennial
clock synchronized to the census. Article I, Section 2 mandates that representatives
``be apportioned among the several States which may be included within this Union,
according to their respective Numbers.'' Since 1941, the method of equal proportions
(Huntington-Hill) has been the statutory apportionment method under 2 U.S.C.\ \S 2a.
The 2030 census will trigger the 85th application of Huntington-Hill, and the results
will determine which states gain or lose congressional seats---and therefore which states
must undergo substantive redistricting (as opposed to minor adjustments).

For the \textsc{bisect} platform, the projected 2030 reapportionment produces two
categories of planning input:

\begin{enumerate}
\item \textbf{New $k$-values}: States gaining or losing seats require updated
  \textsc{bisect} configurations with the new $k$ parameter. The new $k$ affects
  the bisection tree structure (via the \texttt{prime-factor} or \texttt{standard-bisect}
  structure choice), the resolution classification (via the F.3 $k/n > 0.05$ rule),
  and the convergence threshold (via the $T=600$ convergence search calibration).

\item \textbf{Stability forecasts}: States with stable $k$-values but large population
  shifts (Sun Belt growth in Texas, Florida, Arizona) may experience substantial
  redistricting changes even without seat changes, because the population distribution
  across counties changes in ways that alter the optimal county groupings.
\end{enumerate}

This paper provides the projected 2030 $k$-values (Section~\ref{sec:apportionment}),
the cross-census stability analysis (Section~\ref{sec:stability}), and the Texas
$k=41$ prime-factor fallback analysis (Section~\ref{sec:texas}).

\subsection{Data Sources}

The primary data source for the reapportionment projections is the Census Bureau's
2023 National Population Projections, main series (medium variant), which provides
state-level population projections at annual intervals through 2060 \citep{censusProjections2023}.
We extract projected state populations for April 1, 2030 (census day), including
resident population plus the overseas military adjustment that is included in the
apportionment count but not the resident count.

Overseas military and federal civilian employees abroad are allocated to their home
states using the 2020 cycle's allocation formula as a proxy, updated for 2030 using
2022 Defense Manpower Data Center counts by state-of-legal-residence.

\subsection{Relationship to J-Series Apportionment Papers}

The J-series papers (J.0--J.5) provide the mathematical foundations for Huntington-Hill
and alternative apportionment methods. Q.4 applies those foundations to the specific
context of 2030 projections and algorithmic redistricting implications. The
\textsc{bisect-apportion} crate, documented in J.0, implements the Huntington-Hill
algorithm and is used directly in this paper's computations. All seat allocations
in this paper can be reproduced with:

\begin{verbatim}
bisect apportion --year 2030 --input projections_2030.csv \
  --method huntington-hill
\end{verbatim}
